\documentclass{article}
\usepackage[left=2cm, right=2cm, top=2cm, bottom=2cm]{geometry}
\usepackage{graphicx}
\usepackage{amsmath}
\usepackage{amssymb}
\usepackage{amsthm}
\usepackage{fancyhdr}
\usepackage{verbatim}
\usepackage{listings}
\usepackage{xcolor}
\usepackage{pdfpages}
\usepackage{cancel}
\usepackage{physics}
\usepackage{esint}
\usepackage{braket}

\lstset{
    language=C++,
    basicstyle=\ttfamily\footnotesize,
    keywordstyle=\color{blue},
    commentstyle=\color{green},
    stringstyle=\color{red},
    numbers=left,
    numberstyle=\tiny\color{gray},
    frame=single,
    breaklines=true,
    captionpos=b,
}


\title{MTH 446: Lecture \# 18}
\author{Cliff Sun}

\newtheorem{theorem}{Theorem}[section]
\newtheorem{lemma}[theorem]{Lemma}
\newtheorem{definition}[theorem]{Definition}
\newtheorem{conjecture}[theorem]{Conjecture}
\newtheorem{proposition}[theorem]{Proposition}
\newtheorem{corollary}[theorem]{Corollary}
\newtheorem{one minute paper}[theorem]{One Minute Paper}

\pagestyle{fancy}
\lhead{\textbf{\thepage}\ \ \nouppercase{\rightmark}}
\chead{MTH 446: Lecture \# 18}
\rhead{Cliff Sun}

\begin{document}

\maketitle

\section*{Lecture Span}
\begin{enumerate}
    \item Hyperbolic complex functions
\end{enumerate}

\section*{Complex trig}

Note that $\sinh$ and $\cosh$ are periodic functions with period $2\pi i$. We can then define 

\begin{enumerate}
    \item $\tanh = \sinh/\cosh$
    \item $\coth = \cosh/\sinh$. 
    \item $\sech = 1/\cosh$
    \item $\csch = 1/\sinh$
\end{enumerate}

Here, such functions are defined for $z$ values such that the denominator is not zero. We can also define derivatives of the hyperbolic trigs. 

\section*{Inverse trig functions}

\begin{equation}
    w = \arcsin(z) \iff z = \sin(w)
\end{equation}

Note we can rewrite 

\begin{equation}
    z = \frac{\exp(iw) + \exp(-iw)}{2i}
\end{equation}

Then, we can solve for $w$ in terms of z, as (letting $t = \exp(i\omega)$)

\begin{equation}
    t^2 - 2izt - 1 = 0
\end{equation}

Then, 

\begin{equation}
    t = \frac{2iz \pm [(2iz)^2 + 4]^{1/2}}{2}
\end{equation}

Then, you obtain 

\begin{equation}
    w = \arcsin(z) = -i\log(iz + (i - z^2)^{1/2})
\end{equation}

Defining the domain of definition is difficult. We will just assume that $z_0$ is defined for $\arcsin$. 

\subsection*{Example}

Calculate $\arcsin(i)$. Then, using the formula, we obtain 

\begin{equation}
    \arcsin = -i\log(-1 \pm \sqrt{2})
\end{equation}

Then, 

\begin{equation}
    \log(-1 \pm \sqrt{2}) = \begin{cases}
        \ln(-1 + \sqrt{2}) + 2\pi n i \\
        \ln(1 + \sqrt{2}) + (\pi + 2\pi n)i
    \end{cases}
\end{equation}

Note, that we can rewrite this guy as 

\begin{equation}
    \arcsin(i) = \pi n + (-1)^{n} i \ln(\sqrt{2} + 1)
\end{equation}

\section*{Derivative of $\arcsin$}

We use the single-valued branch of $\log$. Then 

\begin{equation}
    \arcsin' = \frac{1}{(1-z^2)^{1/2}}
\end{equation}

\section*{Complex-valued functions of real variable}

An interval $I \subseteq \mathbb{R}$ is \underline{non-degenerate} if it contains two or more points. Therefore, $I$ contains infinitely many points. Consider functions $w$ on $I$, a non-degenerate interval. Then, 

\begin{equation}
    w(t) = u(t) + iv(t)
\end{equation}

For derivatives, we have 

\begin{equation}
    w' = u' + iv'
\end{equation}

Also, 

\begin{enumerate}
    \item $D_t(z_0 w) = z_0 w'$
    \item $D_t(\exp(z_0 t)) = z_0 \exp(z_0 t)$
\end{enumerate}



\end{document}